\documentclass[conference]{IEEEtran}
\IEEEoverridecommandlockouts

\usepackage{cite}
\usepackage{amsmath,amssymb,amsfonts}
\usepackage{algorithmic}
\usepackage{graphicx}
\usepackage{textcomp}
\usepackage{xcolor}
\usepackage{booktabs}
\usepackage{array}

\begin{document}

\title{Improved Priority Exchange Server: Scientific Contextualization\\
{\footnotesize Milestone 2 --- Special Emphasis: Real-Time Systems, HSHL}}

\author{\IEEEauthorblockN{Sumon Boiddo}
\IEEEauthorblockA{\textit{BEng in Electronic Engineering} \\
\textit{Hochschule Hamm-Lippstadt}\\
Matriculation Number: 2220054 \\
Group B, Team 6}}

\maketitle

\begin{abstract}
This document is the Milestone~2 (Scientific Contextualization) for the
seminar topic \emph{Improved Priority Exchange Server} (IPE). It places IPE
within the landscape of aperiodic servers, compares it with its neighbours
(PE, DPE, EDL, and TBS) through a comparison table and detailed pairwise
comparisons, analyses its scientific and engineering trade-offs, discusses
its consequences for embedded systems, and explains which related work is
excluded and why. The primary reference is Buttazzo, \textit{Hard Real-Time
Computing Systems}~\cite{buttazzo}, Chapter~6.
\end{abstract}

% ============================================================
\section{Overview of the Aperiodic Server Landscape}
% ============================================================

\begin{itemize}
  \item Aperiodic servers are used to serve soft aperiodic tasks together with
        hard periodic tasks. They are grouped into two families, based on the
        scheduler used for the periodic tasks~\cite{buttazzo}.
  \item The \emph{fixed-priority} (Rate Monotonic, RM) family uses fixed
        priorities. It includes the Polling Server (PS), the Deferrable Server
        (DS), the Priority Exchange (PE) server, and the Sporadic Server (SS).
  \item The \emph{dynamic-priority} (Earliest Deadline First, EDF) family uses
        deadlines to decide priority. It includes the Dynamic Priority
        Exchange (DPE) server, the Dynamic Sporadic Server (DSS), the Total
        Bandwidth Server (TBS), the Earliest Deadline Late (EDL) server, and
        the Improved Priority Exchange (IPE) server.
  \item IPE belongs to the dynamic-priority (EDF) family. Under EDF the
        periodic utilization bound is $100\%$, compared with about $69\%$ for
        RM with many tasks, which allows larger servers and better aperiodic
        responsiveness~\cite{buttazzo}.
  \item IPE comes from a clear family tree: the fixed-priority PE server was
        first extended to EDF as the DPE server, and DPE was then improved into
        IPE by using the precomputed idle times of the EDL server. The lineage
        is therefore PE $\rightarrow$ DPE $\rightarrow$ IPE~\cite{spuri1996}.
\end{itemize}

% ============================================================
\section{Comparison of Aperiodic Server Approaches}
% ============================================================

Table~\ref{tab:comparison} compares IPE with its four closest neighbours ---
PE, DPE, EDL, and TBS --- across three aspects: the scheduler each one uses,
its responsiveness (how quickly it serves aperiodic requests), and its main
limitation. The table highlights three points. First, only PE uses fixed
priority, while all the other servers use EDF. Second, EDL is optimal and IPE
is nearly optimal, so these two give the best responsiveness. Third, each
server pays a different price: EDL pays in heavy run-time computation, IPE pays
in memory, and TBS accepts a small loss of optimality in exchange for
simplicity.

\begin{table*}[t]
\centering
\caption{Comparison of Aperiodic Server Approaches}
\label{tab:comparison}
\renewcommand{\arraystretch}{1.3}
\begin{tabular}{|p{3.0cm}|p{2.0cm}|p{3.0cm}|p{5.5cm}|}
\hline
\textbf{Server} & \textbf{Scheduler} & \textbf{Responsiveness (speed)} &
\textbf{Main limitation} \\
\hline
PE (Priority Exchange) & RM (fixed) & Moderate &
Fixed-priority only; less responsive than the EDF servers \\
\hline
DPE (Dynamic Priority Exchange) & EDF & Good &
Not optimal; replenishment less efficient than IPE \\
\hline
EDL (Earliest Deadline Late) & EDF & \textbf{Optimal (best)} &
Heavy run-time calculation; too expensive to use in practice \\
\hline
TBS (Total Bandwidth Server) & EDF & Very good &
Not optimal; but very simple and lightweight \\
\hline
\textbf{IPE (Improved Priority Exchange)} & \textbf{EDF} &
\textbf{Nearly optimal} &
High memory cost (the stored tables become very large when the periods are
non-harmonic) \\
\hline
\end{tabular}
\end{table*}

% ============================================================
\section{Detailed Comparison: IPE vs. Neighbouring Servers}
% ============================================================

\noindent\textbf{IPE vs. PE.}
PE is the fixed-priority (RM) ancestor of IPE. IPE moves this idea to EDF and
adds the precomputed EDL idle times. As a result, IPE serves aperiodic tasks
much faster (nearly optimal), whereas PE, being fixed-priority, is less
responsive~\cite{spuri1996}.

\noindent\textbf{IPE vs. DPE.}
DPE is IPE's direct base; both use EDF and the same priority-exchange
mechanism. The improvement lies in the replenishment policy: DPE uses a
periodic refill, while IPE refills using the EDL idle times precomputed
offline (the $E$ and $D$ arrays) and always runs at the highest priority. This
makes IPE more efficient and more responsive than DPE~\cite{spuri1996}.

\noindent\textbf{IPE vs. EDL.}
EDL is the optimal server (minimum aperiodic response time), and IPE nearly
matches it. The key difference is \emph{when} the computation is performed:
EDL recomputes the schedule at run time (online), which is expensive, while
IPE precomputes the same idle-time information offline, which is far cheaper at
run time. IPE therefore achieves nearly the same responsiveness as EDL at a
much lower run-time cost, paying instead in memory~\cite{buttazzo}.

\noindent\textbf{IPE vs. TBS.}
TBS is another EDF server, but it is much simpler: it assigns each aperiodic
request a deadline and schedules it directly under EDF, with almost no overhead
and no stored tables. IPE is more complex, requiring the precomputed $E$ and
$D$ arrays and the priority-exchange bookkeeping, so it needs more memory. The
trade-off is that TBS is lighter and simpler but not optimal, while IPE is
nearly optimal --- especially under heavy aperiodic load --- at the cost of
offline precomputation and additional memory~\cite{buttazzo}.

% ============================================================
\section{Scientific and Engineering Trade-offs}
% ============================================================

\begin{itemize}
  \item \textbf{Optimality vs. practicality.} The EDL server is optimal but
        impractical, because it recomputes the schedule at run time. IPE gives
        up a very small amount of optimality (it is \emph{nearly} optimal
        instead of optimal) in exchange for becoming practical, since its heavy
        work is done offline. This is the central trade-off that defines IPE.
  \item \textbf{Run-time cost vs. memory cost.} IPE moves the heavy idle-time
        calculation from run time to an offline step. This saves run-time cost
        but must be paid for in memory: the precomputed $E$ and $D$ arrays grow
        with the hyperperiod and become very large when the periods are
        non-harmonic.
  \item \textbf{Performance vs. flexibility.} Because everything is precomputed
        for a fixed, known task set, IPE performs very well on static systems
        but cannot adapt to dynamic systems. If the task set changes at run
        time, the precomputed plan becomes invalid and the guarantees are lost.
        IPE thus trades flexibility for high performance.
  \item \textbf{Near-optimal performance vs. simplicity.} Compared with TBS,
        IPE is nearly optimal, especially under heavy load, but it is more
        complex (precomputed tables plus priority-exchange bookkeeping). TBS is
        much simpler and lighter but not optimal. IPE therefore trades
        simplicity for higher responsiveness.
\end{itemize}

% ============================================================
\section{Embedded Systems Consequences}
% ============================================================

\begin{itemize}
  \item To deploy IPE, the operating system must provide an EDF scheduler, and
        an offline analysis step must compute and store the $E$ and $D$ arrays
        before the system starts running. The heavy idle-time computation is
        therefore a design-time activity, not a run-time one.
  \item IPE requires enough memory to store the $E$ and $D$ tables, whose size
        depends on the hyperperiod $H = \mathrm{lcm}(T_1,\dots,T_n)$. Since
        embedded devices often have limited memory, this is a critical
        practical constraint.
  \item \textbf{Consequence for small devices:} IPE fits well when the periods
        are harmonic (small hyperperiod, small tables, little memory). But when
        the periods are non-harmonic, the tables become very large and may not
        fit in the limited memory of a small device, which makes IPE unsuitable
        there.
  \item \textbf{Overall:} IPE is best suited to static embedded systems with a
        fixed, known task set, harmonic periods, an EDF-capable operating
        system, and a need for fast aperiodic response while guaranteeing all
        hard deadlines. It is unsuitable for dynamic systems (changing task
        set), non-harmonic period sets (memory too large), or
        fixed-priority-only platforms.
\end{itemize}

% ============================================================
\section{Why Some Related Work Is Excluded}
% ============================================================

\begin{itemize}
  \item This briefing focuses only on the servers in IPE's own family (PE, DPE,
        EDL, TBS). A few related approaches are mentioned but not studied in
        depth:
  \item The \emph{Constant Bandwidth Server (CBS)} is not part of the IPE topic
        and is a separate seminar subject; it also follows a different design
        idea (bandwidth reservation, not priority exchange). For these reasons
        it is excluded from deep study here.
  \item The \emph{Dynamic Sporadic Server (DSS)} lies outside the
        priority-exchange lineage of IPE and is referenced only as a
        performance comparison point.
  \item \emph{Multiprocessor scheduling and shared-resource protocols} are
        covered in separate chapters and fall outside the single-processor,
        fully-preemptive scope assumed for IPE.
\end{itemize}

% ============================================================
% References
% ============================================================

\begin{thebibliography}{9}

\bibitem{buttazzo}
G.~C.~Buttazzo,
\textit{Hard Real-Time Computing Systems: Predictable Scheduling Algorithms
and Applications}, 3rd~ed.
Springer, 2011, ch.~6 (Dynamic Priority Servers).

\bibitem{spuri1994}
M.~Spuri and G.~Buttazzo,
``Efficient Aperiodic Service under Earliest Deadline Scheduling,''
in \textit{Proc. IEEE Real-Time Systems Symposium (RTSS)}, 1994.

\bibitem{spuri1996}
M.~Spuri and G.~Buttazzo,
``Scheduling Aperiodic Tasks in Dynamic Priority Systems,''
\textit{Real-Time Systems}, vol.~10, no.~2, March 1996.

\bibitem{lss1987}
J.~Lehoczky, L.~Sha, and J.~Strosnider,
``Enhanced Aperiodic Responsiveness in Hard Real-Time Environments,''
in \textit{Proc. IEEE Real-Time Systems Symposium (RTSS)}, 1987.

\bibitem{cc1989}
H.~Chetto and M.~Chetto,
``Some Results of the Earliest Deadline Scheduling Algorithm,''
\textit{IEEE Trans. Software Engineering}, vol.~15, no.~10, 1989.

\bibitem{ai_usage_m2}
S.~Boiddo,
\textit{AI Usage Protocol, Milestone 2: Improved Priority Exchange Server},
HSHL, 2026. Protocol version 0.2; JSON and BibTeX files submitted with this
milestone.

\end{thebibliography}

\end{document}
